\section{Limitations}\label{sec:limitations}

\paragraph{Label noise.}
We treat NVD vectors as ground truth even though analysts disagree, NVD revises its records, and reporting organizations sometimes score the same CVE differently.
Section~\ref{sec:discussion} shows several errors that a reviewer would accept as alternative readings.
We pin the data snapshot so that our numbers stay fixed as the labels change.

\paragraph{Single split and seed.}
We evaluate on one test year and train each model once, so the bootstrap intervals capture the sampling of test CVEs and nothing about variation from training.
We fine-tune one 3.8B model for one epoch, and its validation accuracy was still rising at the end (Table~\ref{tab:ckpt}), so larger models or longer training may do better.
We did not run large hosted models, and we trained the fine-tuned model without rebalancing the rare tier.

\paragraph{No ensemble.}
We report the disagreement between the two trained systems as an upper bound on what a combination could achieve, and we did not build or test one.

\paragraph{CVSS version.}
New CVEs increasingly carry CVSS v4.0 vectors.
The formula approach carries over, since v4.0 also computes the score from the vector, but we tested nothing on v4.0 data.

\paragraph{Zero-shot sample.}
We evaluate the zero-shot model on \nSub{} random test CVEs rather than the full year, and the intervals in Table~\ref{tab:subset} reflect the smaller sample.

\paragraph{Input length.}
We use only English descriptions and cut them at 2,000 characters, which affects \nTestClipped\% of test CVEs.
